\documentclass[aspectratio=169,10pt]{beamer}
\usetheme{default}
\usecolortheme{default}
\usepackage[utf8]{inputenc}
\usepackage[T1]{fontenc}
\usepackage[spanish]{babel}
\usepackage{booktabs}
\setbeamertemplate{navigation symbols}{}
\setbeamercolor{structure}{fg=blue!65!black}
\title{Cobertura, actualidad y calidad de imágenes urbanas de Mapillary}
\subtitle{Corredor Magdalena del Mar -- San Isidro -- Miraflores}
\author{Equipo DINAMITA: Jeffry Arturo Hilario Quintana, Jeffry Vargas, Fatima Margarita Villon Zarate}
\date{Week 5 -- DS5343 Data Visualization}

\begin{document}
\begin{frame}\titlepage\end{frame}

\begin{frame}{Problema y motivación}
\begin{itemize}
  \item Lima cuenta con imágenes colaborativas a nivel de calle, pero su cobertura y actualidad son desiguales.
  \item Planificadores urbanos y contribuidores de OpenStreetMap necesitan reconocer vacíos de información y datos desactualizados.
  \item El proyecto convierte metadatos de Mapillary en evidencia espacial, temporal y de calidad de captura.
\end{itemize}
\end{frame}

\begin{frame}{Audiencia y propósito}
\textbf{Audiencia:} planificadores urbanos, gestores municipales, investigadores de movilidad y contribuidores de OpenStreetMap.\medskip

\textbf{Propósito:} permitir explorar dónde, cuándo y con qué calidad se han capturado imágenes urbanas en el corredor de estudio, para identificar sectores prioritarios de actualización.
\end{frame}

\begin{frame}{Dataset principal propuesto}
\begin{itemize}
  \item Área: Magdalena del Mar, San Isidro y Miraflores.
  \item Fuente: Mapillary Graph API; recolección reproducible mediante cuadrícula de cajas pequeñas.
  \item Colección actual: 18 celdas, 9,938 imágenes válidas y 1,915 secuencias.
  \item Periodo disponible: abril de 2015 a agosto de 2026.
  \item Tabla de detecciones: etiqueta de objeto y frecuencia por imagen.
  \item Capa de referencia: límites distritales.
\end{itemize}
\end{frame}

\begin{frame}{Muestra piloto y preparación}
\begin{itemize}
  \item La muestra de 10 imágenes validó el esquema antes de la recolección completa.
  \item El dataset final conserva 9,938 imágenes tras eliminar duplicados y 568 puntos fuera de los límites distritales.
  \item Se convierten tiempos Unix a UTC y se agregan detecciones repetidas en una tabla relacionada de 166,475 filas.
\end{itemize}
\end{frame}

\begin{frame}{Preguntas de dominio}
\begin{enumerate}
  \item ¿Cómo varía la densidad de imágenes entre los tres distritos?
  \item ¿Qué sectores contienen imágenes recientes y cuáles requieren actualización?
  \item ¿Cómo se asocian cámara, panorámica y calidad de captura?
  \item ¿Qué categorías de infraestructura detectada predominan por zona?
  \item ¿Qué secuencias concentran mayor cobertura y qué patrones de orientación muestran?
\end{enumerate}
\end{frame}

\begin{frame}{Factibilidad y plan inmediato}
\begin{itemize}
  \item Jeffry Hilario: adquisición, validación y tablas reproducibles.
  \item Jeffry Vargas: cuadrícula, límites distritales y análisis espacial.
  \item Fatima Villon: preguntas, narrativa, diseño e iteración con usuarios.
  \item Riesgos: límites de la API, cobertura desigual y etiquetas repetidas; se mitigan con cuadrícula, registro de consultas y agregación de datos.
\end{itemize}
\end{frame}
\end{document}
